\paragraph{Continuous $k=2$ bounds (post hoc sensitivity).} The finite grid
understates interval widths for $k=2$. A post hoc analysis of the same saved
measurements, with no new draws, removed the grid restriction for the
.90$\to$.95 two-opportunity target. At fixed $e$ the survival ratio is a ratio of
affine functions of $\rho$, so its extremes over $\rho$ occur at feasible
endpoints; the assay constraints give a single interval of feasible $e$, and on
each clipped $\rho$ boundary prevention is nondecreasing in $e$, so the global
extremes are attained at the smallest and largest feasible $e$ (the proof holds
for $0<h_0<h_1<1$ in this family). The solver used 60-digit decimal interval
arithmetic with directed rounding (24,000 confidence sets; zero unresolved
decisions; enclosures of width at most $10^{-58}$), and an independent 100-digit
implementation verified all 47,684 feasible extrema and all 158 empty-set
decisions. The certificate is for optimization given the supplied assay limits,
which are floating-point Clopper--Pearson quantiles widened outward by
$3\times10^{-12}$; the whole statistical pipeline is not formally machine
certified. A documented verification-only repair allowed $10^{-90}$ scalar
comparison error after roundoff near an exact point-valued certificate; the
solver and outputs were unchanged.

